% Auxiliary ordinary marking comparison for the finite A.4 compiler.
\paragraph{Adjoining the remaining involution fixed points.}
Let $T$ be the oriented torus with the actual involution
$\iota:z\mapsto-z$, marked origin $O$, and four
ordered unordered pairs $P_i=\{a_i,-a_i\}$ of
nonfixed points.  Write
$B=\operatorname{Fix}(\iota)\setminus\{O\}$;
this is a set of three distinct points.  Let $M_0$
be the ordinary mapping-class group retaining $O$
and each named pair $P_i$ setwise.  Let $M_3$ retain
these same markings and also retain $B$ as one
unordered block.  The actual forgetting map
$p:M_3\to M_0$ fixes the named involution, and
\[
 p:C_{M_3}(\iota)\xrightarrow{\ \sim\ }
                         C_{M_0}(\iota).
 \tag{IFM.1}
\]

To prove this, use the same branched double cover
$T\to T/\langle\iota\rangle\cong S^2$ on both
sides.  Its quotient markings are a distinguished
branch point $\overline O$, the other three branch
points as an unordered set, and the four ordered
points $\overline P_i$.  Denote that literal
mixed sphere group by $Q$.
The ordinary centralizer comparison of A.1 gives
the commutative diagram of actual quotient maps
\[
 \begin{array}{ccccccccc}
 1&\longrightarrow&\langle\iota\rangle
   &\longrightarrow&C_{M_3}(\iota)
   &\longrightarrow&Q&\longrightarrow&1\\
 &&\Vert&&\big\downarrow p&&\Vert\\
 1&\longrightarrow&\langle\iota\rangle
   &\longrightarrow&C_{M_0}(\iota)
   &\longrightarrow&Q&\longrightarrow&1.
 \end{array}                                           \tag{IFM.2}
\]
For the map onto $Q$, every quotient mapping class
preserving these colors lifts: the double-cover
monodromy sends all four branch meridians to the
nontrivial element of $C_2$ and sends the four
pair-image meridians to zero.  These values are
preserved by every named colored mapping class.
The two lifts differ exactly by $\iota$.
The same centralizer-to-quotient comparison identifies
the kernel with that deck group before and after
the three extra fixed-point markings.  Forgetting
them changes neither the actual cover nor its
quotient mapping class.  The two identity arrows
in (IFM.2) therefore make its middle arrow both
injective and surjective, proving (IFM.1).

There is a useful exact description of ordered
intermediate choices.  Order the three points of
$B$ temporarily.  The action of $C_{M_0}(\iota)$
on $B$ is the homomorphism
\[
 C_{M_0}(\iota)\longrightarrow Q
                  \longrightarrow\mathfrak S(B).
 \tag{IFM.3}
\]
It is onto: adjacent half-twists among these
three branch marks, along arcs avoiding the five
other quotient marks, give the two adjacent
transpositions and lift through the same cover.
Thus six literal braid words provide coset
representatives.  The image of the centralizer
after adjoining $j$ distinct fixed points with
their chosen individual labels is exactly the
subgroup fixing those $j$ choices in (IFM.3).
The successive indices are $3,2,1$.
This assertion is a statement about those actual
subgroups and forgetting maps, proved by the same
lifting argument with the extra colors retained.

In a finite point-forgetting construction the
three successive one-point orbifold kernels must
therefore retain the distinct available cone
types.  One first retains the six representatives
of (IFM.3), and then uses the corresponding cone
type at each ordered stage.  The ordinary maps
that these stages compare are the stabilizer maps
just specified.  When all three points are again
treated as one unordered block, their sole
residual label action is the common
$\mathfrak S(B)$ of (IFM.3).  The later map
forgetting $B$ on the quotient sphere kills this
permutation while retaining $\overline O$ and
the four named pair images.  Equation (IFM.1)
does not by itself assert a finite centralizer
recognition bound; those bounds are supplied
by the separate point-fibre refinements.

\par\medskip\noindent\textbf{Allowing permutation of all four pairs.}
Let $M_0^{\Sigma}$ retain $O$ and the unordered collection
$\{P_1,P_2,P_3,P_4\}$, permitting arbitrary permutations of the four
two-element blocks, and let $M_3^{\Sigma}$ retain in addition the
unordered block $B$.  These are actual mixed mapping-class groups.
If $\mathcal M=\Gamma_{1,[9]}$ is the full nine-label mixed group,
then
\[
 C_{\mathcal M}(\iota)=C_{M_0^{\Sigma}}(\iota).
 \tag{IFM.4}
\]
Indeed the label permutation of an element centralizing $\iota$
commutes with its permutation consisting of one fixed label and four
two-cycles.  It therefore fixes the unique fixed label $O$ and
permutes exactly the four named two-element orbits.  The reverse
inclusion is the inclusion of these actual subgroups of $\mathcal M$.

Let $Q^{\Sigma}$ be the sphere group retaining the distinguished
branch point $\overline O$, the three other branch marks as one
unordered block, and the four pair images as another unordered block.
The actual forget map $p^{\Sigma}:M_3^{\Sigma}\to M_0^{\Sigma}$ gives
the corresponding exact diagram
\[
 \begin{array}{ccccccccc}
 1&\longrightarrow&\langle\iota\rangle
   &\longrightarrow&C_{M_3^{\Sigma}}(\iota)
   &\longrightarrow&Q^{\Sigma}&\longrightarrow&1\\
 &&\Vert&&\big\downarrow p^{\Sigma}&&\Vert\\
 1&\longrightarrow&\langle\iota\rangle
   &\longrightarrow&C_{M_0^{\Sigma}}(\iota)
   &\longrightarrow&Q^{\Sigma}&\longrightarrow&1.
 \end{array}
 \tag{IFM.5}
\]
Every displayed quotient mapping class still preserves the actual
double-cover monodromy: all four branch meridians have value one in
$C_2$ and every pair-image meridian has value zero.  Thus every such
class lifts, its two lifts differ by $\iota$, and the kernel of either
centralizer quotient is exactly that deck group.  Both outer vertical
maps in (IFM.5) are the identity of the same marked cover.  It follows
that $p^{\Sigma}$ restricts to an isomorphism of the two centralizers.

The passage from the earlier ordered-pair diagram to (IFM.5) is
itself given by the actual inclusions of marking groups.  In either
row its extra quotient is precisely $\mathfrak S_4$: the kernel
fixes each pair block, and adjacent half-twists among the four
pair-image marks, along arcs avoiding every branch mark, lift to
representatives of the adjacent block transpositions.  The branch
$\mathfrak S_3$ of (IFM.3) remains independently realizable by arcs
avoiding the pair images.  Hence the ordered fixed-point intermediate
subgroups retain the same successive indices $3,2,1$, with the four
pair blocks still permitted to move.

Finally the exact ordinary branch map used for the full centralizer
in (FBO.1) is
\[
 C_{\mathcal M}(\iota)
   \longrightarrow Q^{\Sigma}
   \xrightarrow{\operatorname{forget}B}
       \Gamma_{0,\{\overline O;[\overline P_1,\ldots,
                                        \overline P_4]\}}.
 \tag{IFM.6}
\]
Its first arrow is the lower quotient map in (IFM.5), using (IFM.4).
The second is the literal forgetting homomorphism.  The three branch
half-twist representatives become the identity after forgetting their
marks, while the root-block permutation is retained.  Composing the
upper quotient map in (IFM.5) with the same forgetting map gives
exactly (IFM.6) after $p^{\Sigma}$; thus the comparison includes the
whole ordinary centralizer and its four-pair permutation action.
